\section{讨论与后续研究}
\label{sec:5}

本文平台面向分布式算法、计算负载与节点拓扑相对固定的测试场景，支持主动设置一组无线收发参数，进而定量比较端到端训练轮次的通信开销。受控比较能够成立，依赖底层配置的确定性，以及通信对端对帧级元数据的准确记录。这就对硬件提出了明确要求。实验能力的实际覆盖面，主要受制于所选网卡、射频芯片以及系统驱动的开放程度。不同硬件对频宽切换、调制编码等级以及元数据暴露的支持各不相同，部分消费级网卡驱动并未开放全部底层读写接口。测试环境同样划定了适用边界。在特定室内测得的传输表现与耗时分布，深受房间格局、天线摆放以及背景电磁干扰影响，相关结论无法直接外推至室外远距离传输、高速移动或强干扰等工业车载环境。清楚这些使用条件，才能在合理的边界内展开有针对性的实验设计。

当协同训练的客户端从十余台继续增加时，通信组织的效率权衡开始发生变化。在下行阶段，一对多广播展现出极高的复用价值。无论接收节点增加多少，模型文件在空中只需广播一次，通信负载不会随设备数量成倍增加。上行阶段则完全不同。每个客户端必须独立提交各自计算得到的更新文件。现有平台依托模型广播分发、限时发送权轮转与聚合屏障，保证了多节点交替上传时不会互相冲撞。然而，排队等待许可的节点一旦增多，调度轮询的切换间隙、传输过程中的反向确认，以及慢速节点带来的同步拖尾，都会增加整轮训练的时间负担。下行广播节省的时间能否抵消轮转调度与长尾等待的开销？这一得失平衡在更大规模的网络中究竟如何演变，需要进一步的实证测量。如何在多节点集群中更灵活地协调发送机会与反馈窗口，进而保持整轮协作的高效推进，构成了后续优化通信组织的核心问题。

当训练任务发生变化时，模型体积可能从数兆字节延伸至数百兆字节，本地计算耗时也从数秒跨越到数分钟，各个客户端所处的信道质量更往往参差不齐。在如此差异化的场景下，单一的无线配置很难在整轮训练中始终保持全局高效。显式参数配置与全流程比较提供了观察手段，但不同配置在上下行阶段的收益权衡大相径庭。追求高速率的配置能迅速完成下行模型广播，可一旦遇到信道衰落严重的弱链路节点，上行回传容易出现大量丢包与重传，拖垮整轮进度。低阶调制配合高冗余纠错能够稳妥应对恶劣链路，却又会在下行阶段消耗过多时间。如何在模型规模、计算耗时与链路质量各不相同的情况下确定合理的无线配置，需要更系统的方法支撑。推进这一课题，不仅要在多种神经网络结构与复杂的物理传播环境中展开广泛验证，也需要结合具备更高吞吐特性的新一代无线网卡设备开展长期的对比测试。
